% Physical unit 024; accepted Session 02.
\GordanUnitStart{024}{printed page 19}
\GordanTopFolio{19}
est, $U$ aut positivum aut negativum erit, prout:
\[
 {u\over \sin\varphi}\left\{{u\over \sin\varphi}-{u_2\over \sin\varphi_2}\right\}
\]
aut positivum aut negativum est.

Cui convenit,
\[
 {u\over \sin\varphi}
\]
non iterum evanescente, $U$ positivum esse debere, provecto autem ultra $0$, $U$ aut positivum aut negativum esse posse.

\begin{center}\S\ 15.\end{center}

Unde statuimus theorema:
\begin{quote}
\gq{Cum punctum aliquod, lineae geodeticae proxime situm, per lineas geodeticas jungitur cum punctis $A$ et $B$; summa extensionum earum linearum, si $\psi<\psi'$ est, major extensione \gq{$A\ B$}; si $\psi>\psi$ est, aut minor aut major.}
\end{quote}

\bigskip\hrule\bigskip

\begin{center}
{\large CAPUT V.}\\[0.8em]
\S\ 1.
\end{center}

Jam vero lineae omnium brevissimae, cujus definitionem initio superioris capitis proposuimus, tractanda est discussio. Quam ob rem omnes lineas geodeticas punctum $A(l,\omega)$ continentes considerabimus. Quarum omnes deinceps aequationes eliciuntur, vindicando in:
\[
 \psi-\omega=\pm\int_{\varphi_1}^{\varphi}d\varphi\cos l_0{\sqrt{1-e^2+e^2\sin^2l_0\cos^2\varphi}\over 1-\sin^2l_0\cos^2\varphi}
\]
(in qua angulos $\varphi$ et $\varphi_1$ formulis
\[
 \sin l_0\cos\varphi=\sin l\qquad \sin l_0\cos\varphi_1=\sin l_1
\]
definiri constat) constanti $l_0$ omnes numeros a $0$ ad ${\pi\over2}$; definito autem $l_0$, attamen aequatio manet anceps, quum integrale tum positivum, tum negativum sumi possit; ut binae lineae geodeticae eundem tangant circulum.

\begin{center}\S\ 2.\end{center}

Quarum linearum genera sunt quatuor. Puncto $\Pi$ enim in earum qualibet ad propinquum aliquod punctum $(l,\psi)$ provecto, est:
